% ══════════════════════════════════════════════════════════════════════════════
% Lecture 22: Hashing (Summary) and Quick Sort
% ══════════════════════════════════════════════════════════════════════════════

% ══════════════════════════════════════════════════════════════════════════════
\section*{Hash Tables: Summary (Recap)}
\addcontentsline{toc}{subsection}{Hash Tables: Summary (Recap)}
% ══════════════════════════════════════════════════════════════════════════════

\begin{complexity}{Hash Tables}{hashrecap}
Hash tables efficiently implement the following operations:
\begin{itemize}[noitemsep]
  \item \textbf{Insert:} $O(1)$
  \item \textbf{Delete:} $O(1)$
  \item \textbf{Search:} Expected $O(n/m)$ with a good hash function (i.e.\ $O(1)$ when
        $m = \Theta(n)$)
\end{itemize}
Collisions cannot be avoided without having $m \gg n^2$ (Birthday Lemma). Instead, they are
handled using, for example, \textbf{chaining}.
\end{complexity}

% ══════════════════════════════════════════════════════════════════════════════
\section*{Quick Sort}
\addcontentsline{toc}{subsection}{Quick Sort}
% ══════════════════════════════════════════════════════════════════════════════

% ─────────────────────────────────────────────────────────────────────────────
\subsection{Overview}

\begin{itemize}[noitemsep]
  \item The sorting algorithm of choice in many computer systems.
  \item Easy to implement.
  \item Fast in practice (and as we will prove in theory).
  \item Like merge sort, based on the \textbf{divide-and-conquer} paradigm.
\end{itemize}

% ─────────────────────────────────────────────────────────────────────────────
\subsection{Divide-and-Conquer Structure}

To sort the subarray $A[p \ldots r]$:

\begin{itemize}[noitemsep]
  \item \textbf{Divide:} Partition $A[p \ldots r]$ into two (possibly empty) subarrays
        $A[p \ldots q-1]$ and $A[q+1 \ldots r]$, such that every element in the left subarray
        is $\le A[q]$ and every element in the right subarray is $\ge A[q]$. The element
        $A[q]$ is called the \textbf{pivot}.
  \item \textbf{Conquer:} Sort the two subarrays by recursive calls to \textsc{Quicksort}.
  \item \textbf{Combine:} No work needed — the subarrays are sorted in place.
\end{itemize}

\begin{algorithm}[H]
\caption{Quick Sort}
\begin{algorithmic}[1]
\Procedure{Quicksort}{$A, p, r$}
  \If{$p < r$}
    \State $q \gets$ \Call{Partition}{$A, p, r$}
    \State \Call{Quicksort}{$A, p, q-1$}
    \State \Call{Quicksort}{$A, q+1, r$}
  \EndIf
\EndProcedure
\end{algorithmic}
\end{algorithm}

% ─────────────────────────────────────────────────────────────────────────────
\subsection{The Partition Procedure}

\textsc{Partition} always selects the last element $A[r]$ as the pivot and rearranges the
subarray in place.

\begin{algorithm}[H]
\caption{Partition}
\begin{algorithmic}[1]
\Procedure{Partition}{$A, p, r$}
  \State $x \gets A[r]$ \AlgComment{Pivot}
  \State $i \gets p - 1$
  \For{$j = p$ \textbf{to} $r - 1$}
    \If{$A[j] \le x$}
      \State $i \gets i + 1$
      \State exchange $A[i]$ with $A[j]$
    \EndIf
  \EndFor
  \State exchange $A[i+1]$ with $A[r]$ \AlgComment{Place pivot in correct position}
  \State \Return $i + 1$
\EndProcedure
\end{algorithmic}
\end{algorithm}

% ─────────────────────────────────────────────────────────────────────────────
\subsection{Correctness of Partition}

\begin{theorem}{Loop Invariant of Partition}{partitioninv}
At the start of each iteration of the \textbf{for} loop, for any index $k$:
\begin{enumerate}[noitemsep]
  \item If $p \le k \le i$, then $A[k] \le x$ (pivot).
  \item If $i+1 \le k \le j-1$, then $A[k] > x$.
  \item $A[r] = x$ (pivot remains in place until the final swap).
\end{enumerate}

\textbf{Initialization:} Before the loop, both subarrays $A[p \ldots i]$ and
$A[i+1 \ldots j-1]$ are empty, and $A[r] = x$. The invariant holds trivially.

\textbf{Maintenance:} If $A[j] \le x$, increment $i$ and swap $A[i]$ with $A[j]$, then
increment $j$. If $A[j] > x$, only increment $j$.

\textbf{Termination:} When the loop terminates, $j = r$, so all elements are partitioned
into $A[p \ldots i] \le x$, $A[i+1 \ldots r-1] > x$, and $A[r] = x$. The final two lines
swap $A[i+1]$ with $A[r]$, placing the pivot in its correct position $q = i+1$.
\end{theorem}

% ─────────────────────────────────────────────────────────────────────────────
\subsection{Running Time of Partition}

\begin{complexity}{Partition}{partitioncost}
\begin{itemize}[noitemsep]
  \item The \textbf{for} loop runs $n := r - p + 1$ times.
  \item Each iteration takes $\Theta(1)$.
  \item Total: $\Theta(n)$ for an array of length $n$.
  \item The number of comparisons made is $\approx n$.
\end{itemize}
\end{complexity}

% ══════════════════════════════════════════════════════════════════════════════
\section*{Running Time Analysis of Quick Sort}
\addcontentsline{toc}{subsection}{Running Time Analysis of Quick Sort}
% ══════════════════════════════════════════════════════════════════════════════

% ─────────────────────────────────────────────────────────────────────────────
\subsection{Worst Case}

The worst case occurs when \textsc{Partition} always produces a maximally unbalanced split:
one subarray of size $n-1$ and one of size $0$. This happens when the input is already sorted
(ascending or descending) and the pivot is always the smallest or largest element.

The recursion tree has $n$ levels each costing $\Theta(n), \Theta(n-1), \ldots, \Theta(1)$:
\[
  T(n) = T(n-1) + T(0) + \Theta(n) = T(n-1) + \Theta(n)
  \;\Rightarrow\; T(n) = \Theta(n^2)
\]

% ─────────────────────────────────────────────────────────────────────────────
\subsection{Best Case}

The best case occurs when \textsc{Partition} always produces a perfectly balanced split,
with both subarrays of size $n/2$:
\[
  T(n) = 2T(n/2) + \Theta(n) = \Theta(n \log n)
\]
by the Master Theorem.

% ─────────────────────────────────────────────────────────────────────────────
\subsection{Average Case — Intuition}

\begin{remark}{Unbalanced splits still give $\Theta(n \log n)$}{avgcase}
Even a constant-ratio unbalanced split (e.g.\ always $9$-to-$1$) gives:
\[
  T(n) = T\!\left(\tfrac{9n}{10}\right) + T\!\left(\tfrac{n}{10}\right) + \Theta(n) = \Theta(n \log n)
\]
In practice, splits alternate between good and bad. Alternating best-case and worst-case splits
yields the same asymptotic running time $\Theta(n \log n)$ as always splitting evenly, because
the total work at each level of the tree remains $\Theta(n)$.
\end{remark}

% ══════════════════════════════════════════════════════════════════════════════
\section*{Randomized Quick Sort}
\addcontentsline{toc}{subsection}{Randomized Quick Sort}
% ══════════════════════════════════════════════════════════════════════════════

% ─────────────────────────────────────────────────────────────────────────────
\subsection{Motivation and Algorithm}

The deterministic worst case occurs on sorted inputs — exactly the inputs that arise most
often in practice. To eliminate this vulnerability:

\begin{itemize}[noitemsep]
  \item Instead of always using $A[r]$ as the pivot, \textbf{randomly select} a pivot from
        $A[p \ldots r]$ and swap it with $A[r]$ before calling \textsc{Partition}.
  \item This guarantees good expected performance \textbf{for any input}, not just random ones.
\end{itemize}

\begin{remark}{Key distinction}{randvsdeter}
There is a huge difference between:
\begin{itemize}[noitemsep]
  \item \emph{Expected running time over all inputs} (average-case of deterministic algorithm):
        depends on input distribution, can be manipulated by an adversary.
  \item \emph{Expected running time for any fixed input} (randomized algorithm): the randomness
        is internal; no adversary can force worst-case behaviour.
\end{itemize}
\end{remark}

\begin{algorithm}[H]
\caption{Randomized Quick Sort}
\begin{algorithmic}[1]
\Procedure{Randomized-Partition}{$A, p, r$}
  \State $i \gets \text{\textsc{Random}}(p, r)$ \AlgComment{Pick a random index in $[p, r]$}
  \State exchange $A[r]$ with $A[i]$
  \State \Return \Call{Partition}{$A, p, r$}
\EndProcedure
\Procedure{Randomized-Quicksort}{$A, p, r$}
  \If{$p < r$}
    \State $q \gets$ \Call{Randomized-Partition}{$A, p, r$}
    \State \Call{Randomized-Quicksort}{$A, p, q-1$}
    \State \Call{Randomized-Quicksort}{$A, q+1, r$}
  \EndIf
\EndProcedure
\end{algorithmic}
\end{algorithm}

% ─────────────────────────────────────────────────────────────────────────────
\subsection{Expected Running Time Analysis}

\begin{theorem}{Expected Running Time of Randomized Quicksort}{randqssort}
Randomized \textsc{Quicksort} has expected running time $O(n \log n)$ for \emph{any} input of
size $n$.

\textbf{Proof.}

\textbf{Setup:} The dominant cost is partitioning. Each call to \textsc{Partition} costs
$O(1)$ plus the number of comparisons in the for loop. Since each element is chosen as pivot
at most once, \textsc{Partition} is called at most $n$ times. Let $X$ be the total number of
comparisons across all calls; then the total work is $O(n + X)$.

\textbf{Counting comparisons with indicator variables:} Rename the elements $z_1 < z_2 <
\cdots < z_n$ (sorted order). Define the set $Z_{ij} = \{z_i, z_{i+1}, \ldots, z_j\}$ and
the indicator variable:
\[
  X_{ij} = I\{z_i \text{ is compared to } z_j\}
\]
Since each pair is compared at most once (only when one of them is the pivot), the total
number of comparisons is:
\[
  X = \sum_{i=1}^{n-1} \sum_{j=i+1}^{n} X_{ij}
\]

\textbf{Applying linearity of expectation:}
\[
  \mathbf{E}[X]
  = \sum_{i=1}^{n-1} \sum_{j=i+1}^{n} \mathbf{E}[X_{ij}]
  = \sum_{i=1}^{n-1} \sum_{j=i+1}^{n} \Pr[z_i \text{ is compared to } z_j]
\]

\textbf{Computing the probability:} $z_i$ and $z_j$ are compared if and only if the first
element chosen as pivot from $Z_{ij}$ is either $z_i$ or $z_j$. If any element strictly
between them is chosen first, they are separated into different subarrays and never compared.
Since the pivot is chosen uniformly at random from $Z_{ij}$, which has $j - i + 1$ elements:
\[
  \Pr[z_i \text{ is compared to } z_j] = \frac{2}{j - i + 1}
\]

\textbf{Summing up:} Setting $k = j - i$:
\[
  \mathbf{E}[X]
  = \sum_{i=1}^{n-1} \sum_{j=i+1}^{n} \frac{2}{j-i+1}
  = \sum_{i=1}^{n-1} \sum_{k=1}^{n-i} \frac{2}{k+1}
  < \sum_{i=1}^{n-1} \sum_{k=1}^{n} \frac{2}{k}
  = \sum_{i=1}^{n-1} O(\log n)
  = O(n \log n)
\]
Therefore the expected running time is $O(n + X) = O(n \log n)$.
\end{theorem}

% ─────────────────────────────────────────────────────────────────────────────
\subsection{Summary of Quick Sort}

\begin{tcolorbox}[colback=lightblue, colframe=keyblue, fonttitle=\bfseries,
                  title={Key Points}]
\begin{itemize}[noitemsep]
  \item \textbf{Quick Sort:} Divide-and-conquer sorting algorithm. Pivot partitions the array
        in place; no merge step needed.
  \item \textbf{Partition:} Runs in $\Theta(n)$. Loop invariant guarantees correct placement
        of pivot.
  \item \textbf{Worst case:} $\Theta(n^2)$ — occurs when partitioning is always maximally
        unbalanced (e.g.\ already sorted input with last-element pivot).
  \item \textbf{Best case:} $\Theta(n \log n)$ — perfectly balanced splits every time.
  \item \textbf{Average case intuition:} Even a constant-ratio split (e.g.\ 9-to-1) gives
        $\Theta(n \log n)$. Mixing good and bad splits does not change the asymptotic bound.
  \item \textbf{Randomized Quicksort:} Pick the pivot uniformly at random from the subarray.
        Guarantees expected $O(n \log n)$ for \emph{any} input, foiling adversarial inputs.
  \item \textbf{Key proof idea:} $z_i$ and $z_j$ are compared iff one of them is the first
        pivot chosen from $Z_{ij}$, giving $\Pr[z_i \text{ vs.\ } z_j] = \frac{2}{j-i+1}$.
        Summing by linearity of expectation yields $\mathbf{E}[X] = O(n \log n)$.
  \item \textbf{Practical advantages:} In-place, small constants, very efficient in practice.
\end{itemize}
\end{tcolorbox}